\documentclass{article}
\usepackage{amsmath,amsthm}
\newtheorem{lemma}{Lemma}
\newtheorem{theorem}{Theorem}

\title{A uniform estimate for geometric response profiles}
\author{}
\date{}

\begin{document}
\maketitle

\section{Setup}

Let \(I=[0,1)\) and \(a_n(x)=x^n\). Recall the finite-cover principle in
the form used below: every open cover of a compact interval admits a finite
subcover.

\begin{lemma}\label{lem:pointwise-neighbourhood}
For every \(x\in I\) and \(\varepsilon>0\), there are \(N_x\) and a
relative open neighbourhood \(U_x\) of \(x\) such that
\[
 y^n<\varepsilon \qquad(y\in U_x,\ n\geq N_x).
\]
\end{lemma}

\begin{proof}
Choose \(N_x\) with \(x^{N_x}<\varepsilon\), which is possible because
\(x^n\to0\) for every fixed \(x<1\). Continuity of \(y^{N_x}\) gives
\(U_x\), and \(y^n\leq y^{N_x}\) for \(n\geq N_x\).
\end{proof}

\section{Uniform estimate}

\begin{theorem}\label{thm:uniform-decay}
For every \(\varepsilon>0\) there is \(N\) such that
\[
 x^n<\varepsilon
 \qquad\text{for all }x\in I\text{ and all }n\geq N.
\]
\end{theorem}

\begin{proof}
Fix \(\varepsilon>0\) and use Lemma~\ref{lem:pointwise-neighbourhood}.
The sets \(U_x\) form an open cover of the bounded interval \(I\).
By the finite-cover principle, choose
\(U_{x_1},\ldots,U_{x_m}\) covering \(I\). With
\(N=\max_j N_{x_j}\), any \(x\in I\) lies in some \(U_{x_j}\), and for
\(n\geq N\),
\[
 x^n\leq x^{N_{x_j}}<\varepsilon.
\]
\end{proof}

\end{document}
